\documentclass[12pt]{article}

\usepackage[a4paper, margin=2.5cm]{geometry}
\usepackage{mathptmx}
\usepackage{setspace}
\usepackage{parskip}
\usepackage{titlesec}
\usepackage{enumitem}
\usepackage[T1]{fontenc}
\usepackage[utf8]{inputenc}
\usepackage{hyperref}

\setstretch{1.5}
\setlength{\parskip}{4pt}

\titlespacing*{\section}{0pt}{8pt}{2pt}
\titleformat{\section}{\normalfont\bfseries\normalsize}{}{0pt}{}

\hypersetup{colorlinks=false, hidelinks}

\begin{document}

\begin{center}
    {\large\bfseries Individual Contribution Report}\\[4pt]
    CMPE 321 — Project 3: Modular DBMS Engine \quad|\quad Boğaziçi University, Spring 2026
\end{center}
\hrule\vspace{4pt}
\noindent\textbf{Name:} Enes \"{O}zt\"{u}rk \hfill \textbf{Student ID:} 2022400042 \hfill \textbf{Group:} 28
\vspace{4pt}\hrule

\section{Tasks \& Contributions}

I was responsible for Phase 4 (experimental evaluation and deliverables) and the shared document structure used across the project.

\begin{itemize}[leftmargin=*, itemsep=1pt, topsep=2pt]
    \item \textbf{Phase 4 — Workload Generator.} Implemented \texttt{workload\_generator.py} with four modes (\texttt{sequential}, \texttt{random}, \texttt{range}, \texttt{mixed}), a fixed \texttt{random.seed(42)} for reproducibility, and a \texttt{create type} header followed by $N$ insert and $Q$ query lines.

    \item \textbf{Phase 4 — Experiments.} Implemented \texttt{run\_experiments.py} running three experiments: (1) LRU vs.\ MRU hit rate and I/O under sequential and random workloads; (2) \texttt{heap\_scan} vs.\ \texttt{hash\_index} vs.\ \texttt{bplus\_tree} query cost; (3) buffer pool size sensitivity (4, 8, 16, 32, 64 frames). Collected results and wrote the analysis tables in \texttt{report.md}.

    \item \textbf{Deliverables.} Authored \texttt{README.md}, \texttt{record.txt} (exact reproduction commands), and the \LaTeX{} individual contribution format used by both teammates.

    \item \textbf{Video.} Co-recorded the project demonstration video covering experiment methodology and live reproduction.
\end{itemize}

\section{Collaboration \& Challenges}

Berkay owned Phases 0–3 while I owned Phase 4. We coordinated through \texttt{IMPLEMENTATION\_GUIDE.md}, which either of us updated immediately after any design decision. The main challenge was that buffer counters accumulate across a session, so experiment scripts needed a \texttt{stats reset} between runs to isolate per-workload numbers. Understanding this required reading Berkay's counter-reset implementation, which highlighted the value of our shared working contract.

\section{Self-Assessment}

Phase 4 required experimental design thinking (controlling variables, avoiding confounds) rather than low-level data structures, which was a valuable shift in perspective. I became more comfortable writing reproducible benchmark harnesses. I would start experiments earlier in future projects to allow time to iterate on workload parameters.

\section{Use of AI Tools}

During the project I used an AI tool mainly to read up on things I was
unfamiliar with. Before designing the experiments I asked it to explain what
sequential flooding means in the context of buffer replacement, which helped me
set up the LRU vs.\ MRU comparison more meaningfully. I also used it a couple
of times to understand how \texttt{argparse} works and to get the \LaTeX{}
table syntax right for the report. The experiment scripts and all written
content are my own work.

\end{document}
